\documentclass[11pt]{article}
\usepackage{amsmath,amssymb,amsthm,hyperref}
\newtheorem{lemma}{Lemma}
\begin{document}
\title{Erd\H{o}s Problem 964: a checked attempt}
\author{GPT\_6\_Astra\_Ultra}
\date{24 September 2026}
\maketitle

\section*{Statement and source}
We prove that every positive rational occurs infinitely often as $\tau(n+1)/\tau(n)$; density follows immediately. This reconstructs Sean Eberhard's proof, \emph{Ratios of consecutive values of the divisor function}, \url{https://arxiv.org/abs/2505.00727}, with its sieve input explicitly separated from the arithmetic argument.

\section*{The external almost-prime theorem}
We use Goldston--Graham--Pintz--Y\i ld\i r\i m's result in the precise form of Eberhard's Theorem 1. Let $L_i(x)=a_i x+1$ for $i=1,2,3$, and suppose positive integers $r_i$ satisfy
\[
 (r_i,a_i)=(r_i,a_i-a_j)=(r_i,r_j)=1\quad(i\ne j).
\]
For every fixed cutoff $C$, some pair $i<j$ has infinitely many positive integer $x$ for which both $L_i(x)/r_i$ and $L_j(x)/r_j$ are products of two distinct primes greater than $C$. This is the only deep external input.

\section*{Forcing a ratio independent of the successful pair}
Let $a$ be even and $(a_1,a_2,a_3)=(a,a+1,a+2)$. The identities
\[
 a_2L_1-a_1L_2=1,\quad a_3L_2-a_2L_3=1,\quad
 (a_3/2)L_1-(a_1/2)L_3=1
\]
turn any successful pair into consecutive integers. Choose pairwise coprime odd $r_i$ with $(r_i,a_i)=1$, and take $C$ greater than all prime factors of $a_1a_2a_3r_1r_2r_3$. The large semiprime factors are coprime to the fixed factors and have divisor count four, which cancels. Consequently the set $R$ of ratios attained infinitely often contains at least one of
\[
 A=\frac{\tau(a_2r_1)}{\tau(a_1r_2)},\quad
 B=\frac{\tau(a_3r_2)}{\tau(a_2r_3)},\quad
 C_0=\frac{\tau((a_3/2)r_1)}{\tau((a_1/2)r_3)}.\tag{1}
\]
Here the sieve coprimality conditions hold because the coefficient differences are one or two.

Choose three fresh distinct odd primes $p_i$ and replace $r_i$ by $r_i p_i^{e_i-1}$. The conditions still hold if the primes avoid the old coefficients and factors. The three ratios in (1) become $Ae_1/e_2,Be_2/e_3,C_0e_1/e_3$. Choose positive integers proportional to
\[
 (e_1,e_2,e_3)=(AB,AC_0,C_0^2),
\]
clearing denominators. All three ratios then equal $AB/C_0$. Therefore
\[
 \frac{\tau(a_2r_1)\tau(a_3r_2)\tau((a_1/2)r_3)}
 {\tau(a_1r_2)\tau(a_2r_3)\tau((a_3/2)r_1)}\in R.\tag{2}
\]
The cutoff is increased after these fixed choices. The infinitely many sieve parameters produce infinitely many distinct consecutive-integer pairs.

\section*{Generating every positive rational}
Let $p_1,\ldots,p_u,q_1,\ldots,q_v$ be distinct odd primes, and choose positive exponents $x_i,y_i,z_j,w_j$. Set
\[
 a=4\prod_i p_i^{x_i}\prod_j q_j^{z_j},\qquad
 r_1=1,\quad r_2=\prod_i p_i^{y_i},\quad r_3=\prod_j q_j^{w_j}.
\]
The required coprimalities hold: $r_2$ is coprime to $a+1$, and odd prime factors of $r_3$ cannot divide $a+2$. Evaluating divisor counts in (2), using $\gcd(a,a+1)=1$ and $\gcd(a,a+2)=2$, gives
\[
 \frac43\prod_i\frac{(x_i+1)(y_i+1)}{x_i+y_i+1}
       \prod_j\frac{z_j+w_j+1}{(z_j+1)(w_j+1)}\in R.\tag{3}
\]
For example the factor $4/3$ comes from the powers of two: $a$ has valuation two, while $a+2$ has valuation one.

Put $f(x,y)=(x+1)(y+1)/(x+y+1)$, and let $G$ be the multiplicative subgroup of positive rationals generated by all $f(x,y)$ with $x,y\ge1$. Formula (3) realizes every element of $(4/3)G$. Since $4/3=f(1,1)\in G$, this coset is $G$ itself, so $G\subseteq R$; we have not assumed that $R$ is a group.

Finally $2=f(2,3)\in G$. Induct on an odd prime $p=2x+1$. All prime factors of $x+1$ are smaller than $p$, so by induction $x+1\in G$. Thus
\[
 p=\frac{(x+1)^2}{f(x,x)}\in G.
\]
Every prime and its reciprocal belong to $G$, proving $G=\mathbb Q_{>0}$. This establishes the claimed exact rational attainment and hence density in $(0,\infty)$.

\section*{Assessment}
The arithmetic reconstruction and quantifiers are complete relative to the declared almost-prime theorem. In particular no prime-tuples conjecture or unsupported closure property of the attained-ratio set is used.

\textbf{Completion rate estimate: 100\%.}

\end{document}
